\documentclass[11pt,a4paper]{article}
\usepackage[utf8]{inputenc}
\usepackage[T1]{fontenc}
\usepackage[catalan]{babel}
\usepackage{lmodern}
\usepackage[sfdefault,lining]{FiraSans}
\usepackage[a4paper,top=2.3cm,bottom=2.2cm,left=1.8cm,right=1.8cm,headheight=1.3cm,headsep=0.4cm,footskip=1cm]{geometry}
\usepackage{amsmath,amssymb}
\usepackage{icomma}
\usepackage{graphicx}
\usepackage[table]{xcolor}
\usepackage{tikz}
\usepackage{fontawesome5}
\usepackage{ccicons}
\usepackage{fancyhdr}
\usepackage{lastpage}
\usepackage{tasks}
\usepackage{needspace}
\usepackage[most]{tcolorbox}
\usepackage{array}
\usepackage[hidelinks,pdftitle={Solucions de la fitxa de la unitat 3: Nombres enters. 1r d'ESO},pdfsubject={Matemàtiques 1r ESO},pdfkeywords={nombres enters, recta numèrica, valor absolut, oposat, suma i resta, regla dels signes, operacions combinades, problemes}]{hyperref}

\definecolor{navy}{HTML}{1B2748}
\definecolor{accent}{HTML}{E8603A}
\definecolor{teal}{HTML}{1F9E90}
\definecolor{soft}{HTML}{EEF1F8}
\definecolor{muted}{HTML}{5D6685}

\setlength{\parindent}{0pt}
\setlength{\parskip}{3pt}
\settasks{label=\alph*), label-format=\bfseries\color{navy}, label-width=1.3em, item-indent=1.8em, column-sep=0.6em, after-item-skip=2pt, before-skip=3pt, after-skip=4pt}

% ---------- capçalera i peu ----------
\newcommand{\llicencia}{\href{https://creativecommons.org/licenses/by-nc-sa/4.0/deed.ca}{\ccbyncsa}\ \ Llicència Creative Commons Reconeixement-NoComercial-CompartirIgual 4.0 Internacional (CC BY-NC-SA 4.0)}
\pagestyle{fancy}
\fancyhf{}
\fancyhead[L]{\raisebox{-0.35\height}{\includegraphics[height=0.95cm]{logo-gva.png}}}
\fancyhead[C]{\small\color{muted}Solucions de la fitxa · Unitat 3: Nombres enters · 1r d'ESO}
\fancyhead[R]{\raisebox{-0.35\height}{\includegraphics[height=0.95cm]{logo-vermellar.png}}}
\fancyfoot[L]{\scriptsize\color{muted}\llicencia}
\fancyfoot[R]{\small\color{muted}\thepage}
\renewcommand{\headrulewidth}{0.4pt}
\renewcommand{\headrule}{\color{muted!40}\hrule height \headrulewidth}
\fancypagestyle{primera}{\fancyhf{}\renewcommand{\headrulewidth}{0pt}
  \fancyfoot[L]{\scriptsize\color{muted}\llicencia}
  \fancyfoot[R]{\small\color{muted}\thepage}}

% ---------- nivells i exercicis ----------
\newcommand{\estrelles}[1]{{\color{accent}\ifcase#1\or\faStar\faStar[regular]\faStar[regular]\or\faStar\faStar\faStar[regular]\or\faStar\faStar\faStar\fi}}
\newcounter{exer}
\newcommand{\ex}[1]{\par\needspace{4\baselineskip}\medskip\stepcounter{exer}\noindent{\color{navy}\bfseries\theexer.}\ \ {\small\estrelles{#1}}\ \ }

\newcommand{\seccio}[2]{\par\needspace{10\baselineskip}\bigskip
  \begin{tcolorbox}[colback=navy,colframe=navy,coltext=white,arc=2mm,boxrule=0pt,left=3mm,right=3mm,top=1.2mm,bottom=1.2mm]
  {\large\bfseries #1\quad #2}
  \end{tcolorbox}\vspace{-2pt}}

\newcolumntype{C}[1]{>{\centering\arraybackslash}p{#1}}

% quadre de recordatori (només a la fitxa d'exercicis)
\newcommand{\recorda}[1]{\par\smallskip\begin{tcolorbox}[colback=accent!7,colframe=accent!55,arc=2mm,boxrule=0.5pt,left=2.5mm,right=2.5mm,top=1mm,bottom=1mm]\small{\color{accent}\faLightbulb[regular]}\ \textbf{Recorda:}\ #1\end{tcolorbox}}
% recta numèrica entera: \rectabuida{min}{max} i \recta{min}{max}{punts}
\newcommand{\rectabase}[2]{\draw[latex-latex,muted] (#1-0.8,0)--(#2+0.8,0);
  \foreach \i in {#1,...,#2}{\draw[muted] (\i,0.12)--(\i,-0.12) node[below=-1pt,font=\scriptsize,text=muted]{$\i$};}}
\newcommand{\rectabuida}[2]{\begin{tikzpicture}[x=0.78cm,baseline=0]\rectabase{#1}{#2}\end{tikzpicture}}
\newcommand{\recta}[3]{\begin{tikzpicture}[x=0.78cm,baseline=0]\rectabase{#1}{#2}
  \foreach \p in {#3}{\fill[accent] (\p,0) circle (0.11) node[above=1pt,font=\small\bfseries,text=navy]{$\p$};}\end{tikzpicture}}
\newcommand{\res}[1]{\boxed{#1}}

% =====================================================================
% Capçalera FentMates de les fitxes: un color per a cada curs.
% Els colors són els de la pàgina inicial (LEVELS a assets/js/materials.js);
% el logo ve de _eines/logo/pdf/ (python3 _eines/logo/genera.py).
%
% Ús, al final del preàmbul de la unitat:
%   % =====================================================================
% Capçalera FentMates de les fitxes: un color per a cada curs.
% Els colors són els de la pàgina inicial (LEVELS a assets/js/materials.js);
% el logo ve de _eines/logo/pdf/ (python3 _eines/logo/genera.py).
%
% Ús, al final del preàmbul de la unitat:
%   % =====================================================================
% Capçalera FentMates de les fitxes: un color per a cada curs.
% Els colors són els de la pàgina inicial (LEVELS a assets/js/materials.js);
% el logo ve de _eines/logo/pdf/ (python3 _eines/logo/genera.py).
%
% Ús, al final del preàmbul de la unitat:
%   \input{../comu/fentmates.tex}
%   \fmcurs{1eso}
% i, al començament del document:
%   \capcalera{Matemàtiques}{1r ESO}{Fitxa d'exercicis · Unitat 2}{Divisibilitat}
% Opcional abans de \input: \def\fmneed{14} (espai mínim abans de cada secció).
% =====================================================================
\makeatletter
\def\fm@logodir{../../logo/pdf}
\providecommand{\fmneed}{10}
\def\fm@cursnom{}

% clau del curs, color i color del text damunt del color
\newcommand{\fm@curs}[3]{\@namedef{fm@color@#1}{#2}\@namedef{fm@text@#1}{#3}}
\fm@curs{1eso}{E6B84A}{navy}
\fm@curs{2eso}{3B82F6}{white}
\fm@curs{3eso}{FF7A50}{white}
\fm@curs{4eso-a}{A78BFA}{white}
\fm@curs{4eso-b}{6D28D9}{white}
\fm@curs{1batx-mi}{9B1C1C}{white}
\fm@curs{1batx-ccss}{EF6B6B}{white}
\fm@curs{2batx-mii}{1E6B47}{white}
\fm@curs{2batx-ccss}{2BB3A3}{white}

\newcommand{\fmcurs}[1]{%
  \@ifundefined{fm@color@#1}{\PackageError{fentmates}{Curs desconegut: #1}{}}{}%
  \edef\fm@tmp{\noexpand\definecolor{curs}{HTML}{\@nameuse{fm@color@#1}}%
               \noexpand\colorlet{curstext}{\@nameuse{fm@text@#1}}}\fm@tmp
  \colorlet{cursfosc}{curs!72!black}%
  \def\fm@logo{\fm@logodir/vertical-#1.pdf}%
  \def\fm@marca{\fm@logodir/marca-#1.pdf}}

% ---------- capçalera de la primera pàgina ----------
% #1 matèria (petit), #2 curs, #3 tipus de document i unitat, #4 títol
\newcommand{\capcalera}[4]{%
  \gdef\fm@cursnom{#2}%
  \par\noindent
  \begin{tcolorbox}[colback=curs,colframe=curs,boxrule=0pt,arc=2.2mm,width=3.5cm,height=2.15cm,
      valign=center,left=3mm,right=2mm,top=1mm,bottom=1mm,nobeforeafter,box align=center]
    \color{curstext}{\footnotesize #1}\par\vspace{1pt}{\fontsize{19}{21}\selectfont\bfseries #2\par}%
  \end{tcolorbox}\hspace{5mm}%
  \begin{minipage}[c]{\dimexpr\linewidth-3.5cm-5mm-2.9cm\relax}
    {\color{cursfosc}\small\bfseries\MakeUppercase{#3}}\par\vspace{2pt}
    {\color{navy}\fontsize{20}{23}\selectfont\bfseries #4\par}
  \end{minipage}\hfill
  \begin{minipage}[c]{2.6cm}\raggedleft\includegraphics[width=2.6cm]{\fm@logo}\end{minipage}\par
  \vspace{2pt}}

% ---------- barres de secció del color del curs ----------
\renewcommand{\seccio}[2]{\par\needspace{\fmneed\baselineskip}\bigskip
  \begin{tcolorbox}[colback=curs,colframe=curs,coltext=curstext,arc=2mm,boxrule=0pt,left=3mm,right=3mm,top=1.2mm,bottom=1.2mm]
  {\large\bfseries #1\quad #2}
  \end{tcolorbox}\vspace{-2pt}}

% ---------- capçalera i peu de les pàgines ----------
% primera pàgina: logos institucionals al peu
\fancypagestyle{primera}{\fancyhf{}\renewcommand{\headrulewidth}{0pt}
  \fancyfoot[L]{\raisebox{-0.32\height}{\includegraphics[height=0.78cm]{logo-gva.png}}\hspace{5mm}%
    \raisebox{-0.32\height}{\includegraphics[height=0.78cm]{logo-vermellar.png}}}
  \fancyfoot[R]{\scriptsize\color{muted}\href{https://creativecommons.org/licenses/by-nc-sa/4.0/deed.ca}{\ccbyncsa}\ \ CC BY-NC-SA 4.0 · fentmates.org\qquad\small\thepage}}
% la resta: el símbol a l'esquerra i el curs a la dreta, amb el títol al centre com abans
\fancyhead[L]{\raisebox{-0.3\height}{\includegraphics[height=0.62cm]{\fm@marca}}}
\fancyhead[R]{\small\bfseries\color{cursfosc}\fm@cursnom}
\renewcommand{\headrule}{\color{curs!45}\hrule height \headrulewidth}
\makeatother

%   \fmcurs{1eso}
% i, al començament del document:
%   \capcalera{Matemàtiques}{1r ESO}{Fitxa d'exercicis · Unitat 2}{Divisibilitat}
% Opcional abans de \input: \def\fmneed{14} (espai mínim abans de cada secció).
% =====================================================================
\makeatletter
\def\fm@logodir{../../logo/pdf}
\providecommand{\fmneed}{10}
\def\fm@cursnom{}

% clau del curs, color i color del text damunt del color
\newcommand{\fm@curs}[3]{\@namedef{fm@color@#1}{#2}\@namedef{fm@text@#1}{#3}}
\fm@curs{1eso}{E6B84A}{navy}
\fm@curs{2eso}{3B82F6}{white}
\fm@curs{3eso}{FF7A50}{white}
\fm@curs{4eso-a}{A78BFA}{white}
\fm@curs{4eso-b}{6D28D9}{white}
\fm@curs{1batx-mi}{9B1C1C}{white}
\fm@curs{1batx-ccss}{EF6B6B}{white}
\fm@curs{2batx-mii}{1E6B47}{white}
\fm@curs{2batx-ccss}{2BB3A3}{white}

\newcommand{\fmcurs}[1]{%
  \@ifundefined{fm@color@#1}{\PackageError{fentmates}{Curs desconegut: #1}{}}{}%
  \edef\fm@tmp{\noexpand\definecolor{curs}{HTML}{\@nameuse{fm@color@#1}}%
               \noexpand\colorlet{curstext}{\@nameuse{fm@text@#1}}}\fm@tmp
  \colorlet{cursfosc}{curs!72!black}%
  \def\fm@logo{\fm@logodir/vertical-#1.pdf}%
  \def\fm@marca{\fm@logodir/marca-#1.pdf}}

% ---------- capçalera de la primera pàgina ----------
% #1 matèria (petit), #2 curs, #3 tipus de document i unitat, #4 títol
\newcommand{\capcalera}[4]{%
  \gdef\fm@cursnom{#2}%
  \par\noindent
  \begin{tcolorbox}[colback=curs,colframe=curs,boxrule=0pt,arc=2.2mm,width=3.5cm,height=2.15cm,
      valign=center,left=3mm,right=2mm,top=1mm,bottom=1mm,nobeforeafter,box align=center]
    \color{curstext}{\footnotesize #1}\par\vspace{1pt}{\fontsize{19}{21}\selectfont\bfseries #2\par}%
  \end{tcolorbox}\hspace{5mm}%
  \begin{minipage}[c]{\dimexpr\linewidth-3.5cm-5mm-2.9cm\relax}
    {\color{cursfosc}\small\bfseries\MakeUppercase{#3}}\par\vspace{2pt}
    {\color{navy}\fontsize{20}{23}\selectfont\bfseries #4\par}
  \end{minipage}\hfill
  \begin{minipage}[c]{2.6cm}\raggedleft\includegraphics[width=2.6cm]{\fm@logo}\end{minipage}\par
  \vspace{2pt}}

% ---------- barres de secció del color del curs ----------
\renewcommand{\seccio}[2]{\par\needspace{\fmneed\baselineskip}\bigskip
  \begin{tcolorbox}[colback=curs,colframe=curs,coltext=curstext,arc=2mm,boxrule=0pt,left=3mm,right=3mm,top=1.2mm,bottom=1.2mm]
  {\large\bfseries #1\quad #2}
  \end{tcolorbox}\vspace{-2pt}}

% ---------- capçalera i peu de les pàgines ----------
% primera pàgina: logos institucionals al peu
\fancypagestyle{primera}{\fancyhf{}\renewcommand{\headrulewidth}{0pt}
  \fancyfoot[L]{\raisebox{-0.32\height}{\includegraphics[height=0.78cm]{logo-gva.png}}\hspace{5mm}%
    \raisebox{-0.32\height}{\includegraphics[height=0.78cm]{logo-vermellar.png}}}
  \fancyfoot[R]{\scriptsize\color{muted}\href{https://creativecommons.org/licenses/by-nc-sa/4.0/deed.ca}{\ccbyncsa}\ \ CC BY-NC-SA 4.0 · fentmates.org\qquad\small\thepage}}
% la resta: el símbol a l'esquerra i el curs a la dreta, amb el títol al centre com abans
\fancyhead[L]{\raisebox{-0.3\height}{\includegraphics[height=0.62cm]{\fm@marca}}}
\fancyhead[R]{\small\bfseries\color{cursfosc}\fm@cursnom}
\renewcommand{\headrule}{\color{curs!45}\hrule height \headrulewidth}
\makeatother

%   \fmcurs{1eso}
% i, al començament del document:
%   \capcalera{Matemàtiques}{1r ESO}{Fitxa d'exercicis · Unitat 2}{Divisibilitat}
% Opcional abans de \input: \def\fmneed{14} (espai mínim abans de cada secció).
% =====================================================================
\makeatletter
\def\fm@logodir{../../logo/pdf}
\providecommand{\fmneed}{10}
\def\fm@cursnom{}

% clau del curs, color i color del text damunt del color
\newcommand{\fm@curs}[3]{\@namedef{fm@color@#1}{#2}\@namedef{fm@text@#1}{#3}}
\fm@curs{1eso}{E6B84A}{navy}
\fm@curs{2eso}{3B82F6}{white}
\fm@curs{3eso}{FF7A50}{white}
\fm@curs{4eso-a}{A78BFA}{white}
\fm@curs{4eso-b}{6D28D9}{white}
\fm@curs{1batx-mi}{9B1C1C}{white}
\fm@curs{1batx-ccss}{EF6B6B}{white}
\fm@curs{2batx-mii}{1E6B47}{white}
\fm@curs{2batx-ccss}{2BB3A3}{white}

\newcommand{\fmcurs}[1]{%
  \@ifundefined{fm@color@#1}{\PackageError{fentmates}{Curs desconegut: #1}{}}{}%
  \edef\fm@tmp{\noexpand\definecolor{curs}{HTML}{\@nameuse{fm@color@#1}}%
               \noexpand\colorlet{curstext}{\@nameuse{fm@text@#1}}}\fm@tmp
  \colorlet{cursfosc}{curs!72!black}%
  \def\fm@logo{\fm@logodir/vertical-#1.pdf}%
  \def\fm@marca{\fm@logodir/marca-#1.pdf}}

% ---------- capçalera de la primera pàgina ----------
% #1 matèria (petit), #2 curs, #3 tipus de document i unitat, #4 títol
\newcommand{\capcalera}[4]{%
  \gdef\fm@cursnom{#2}%
  \par\noindent
  \begin{tcolorbox}[colback=curs,colframe=curs,boxrule=0pt,arc=2.2mm,width=3.5cm,height=2.15cm,
      valign=center,left=3mm,right=2mm,top=1mm,bottom=1mm,nobeforeafter,box align=center]
    \color{curstext}{\footnotesize #1}\par\vspace{1pt}{\fontsize{19}{21}\selectfont\bfseries #2\par}%
  \end{tcolorbox}\hspace{5mm}%
  \begin{minipage}[c]{\dimexpr\linewidth-3.5cm-5mm-2.9cm\relax}
    {\color{cursfosc}\small\bfseries\MakeUppercase{#3}}\par\vspace{2pt}
    {\color{navy}\fontsize{20}{23}\selectfont\bfseries #4\par}
  \end{minipage}\hfill
  \begin{minipage}[c]{2.6cm}\raggedleft\includegraphics[width=2.6cm]{\fm@logo}\end{minipage}\par
  \vspace{2pt}}

% ---------- barres de secció del color del curs ----------
\renewcommand{\seccio}[2]{\par\needspace{\fmneed\baselineskip}\bigskip
  \begin{tcolorbox}[colback=curs,colframe=curs,coltext=curstext,arc=2mm,boxrule=0pt,left=3mm,right=3mm,top=1.2mm,bottom=1.2mm]
  {\large\bfseries #1\quad #2}
  \end{tcolorbox}\vspace{-2pt}}

% ---------- capçalera i peu de les pàgines ----------
% primera pàgina: logos institucionals al peu
\fancypagestyle{primera}{\fancyhf{}\renewcommand{\headrulewidth}{0pt}
  \fancyfoot[L]{\raisebox{-0.32\height}{\includegraphics[height=0.78cm]{logo-gva.png}}\hspace{5mm}%
    \raisebox{-0.32\height}{\includegraphics[height=0.78cm]{logo-vermellar.png}}}
  \fancyfoot[R]{\scriptsize\color{muted}\href{https://creativecommons.org/licenses/by-nc-sa/4.0/deed.ca}{\ccbyncsa}\ \ CC BY-NC-SA 4.0 · fentmates.org\qquad\small\thepage}}
% la resta: el símbol a l'esquerra i el curs a la dreta, amb el títol al centre com abans
\fancyhead[L]{\raisebox{-0.3\height}{\includegraphics[height=0.62cm]{\fm@marca}}}
\fancyhead[R]{\small\bfseries\color{cursfosc}\fm@cursnom}
\renewcommand{\headrule}{\color{curs!45}\hrule height \headrulewidth}
\makeatother

\fmcurs{1eso}

\begin{document}
\thispagestyle{primera}

\capcalera{Matemàtiques}{1r ESO}{Solucions de la fitxa · Unitat 3}{Nombres enters}
\vspace{8pt}

\begin{tcolorbox}[colback=soft,colframe=soft,arc=2mm,left=3mm,right=3mm,top=1.5mm,bottom=1.5mm]
\small Intenta resoldre cada exercici abans de mirar-ne la solució. Si no t'ix el mateix resultat, repassa els passos o consulta els apunts interactius de la unitat.
\end{tcolorbox}

% =====================================================
\seccio{1}{Els nombres enters}

\ex{1} \begin{tasks}(4)
  \task $\res{+432}$
  \task $\res{-6}$
  \task $\res{-3}$
  \task $\res{-470}$
  \task $\res{0}$
  \task $\res{+15}$
  \task $\res{-1}$
\end{tasks}

\ex{1} Per exemple: temperatures baix zero ($-5$\,°C), soterranis d'un edifici ($-2$), deutes o saldos negatius ($-30$\,€), profunditats sota el nivell de la mar ($-200$\,m), anys abans de Crist ($-470$), gols en contra\ldots

\ex{1} Hi ha moltes respostes. Per exemple: \ a) un bussejador a 12\,m de profunditat; \ b) la planta 8 d'un edifici; \ c) un submarí a 1500\,m de profunditat; \ d) l'any 2026 d.\,C.; \ e) fa 4\,°C baix zero.

\ex{1} Positius: $12,\ +3,\ 1,\ 100$. \ Negatius: $-7,\ {-25},\ {-1}$. \ Ni positiu ni negatiu: $0$. \ Naturals: $\res{12,\ +3,\ 1,\ 100,\ 0}$.

\ex{2} \begin{tasks}(5)
  \task $4$ i $6$
  \task $-2$ i $0$
  \task $-1$ i $1$
  \task $-11$ i $-9$
  \task $-100$ i $-98$
\end{tasks}

\ex{2} \begin{tasks}(3)
  \task $\res{-3,{-2},{-1},0,1,2}$
  \task $\res{-9,{-8},{-7}}$
  \task $\res{-1,0,1}$
\end{tasks}

\ex{3} Del $-49$ al $-1$ n'hi ha 49, el 0 i de l'1 al 49, 49 més: $49+1+49=\res{99}$. \ Entre $-50$ i $0$: del $-49$ al $-1$, $\res{49}$.


% =====================================================
\seccio{2}{Recta numèrica, ordre i comparació}

\ex{1} \par\smallskip{\centering\recta{-9}{9}{-2,0,4,{-1},8,{-3},{-7},1}\par}

\ex{1} \par\smallskip{\centering\recta{-9}{9}{-7,3,1,{-4},6,{-5},{-2},0}\par}\[\res{-7<-5<-4<-2<0<1<3<6}\]

\ex{1} $\res{8>3>1>0>-1>-4>-7>-9}$

\ex{1} \begin{tasks}(4)
  \task $-6\res{<}6$
  \task $0\res{>}-1$
  \task $10\res{>}-1$
  \task $8\res{>}3$
  \task $-8\res{>}-9$
  \task $1\res{>}-2$
  \task $-15\res{>}-51$
  \task $-100\res{<}1$
\end{tasks}

\ex{2} \begin{tasks}(4)
  \task $-4$ o $-3$
  \task $\res{0}$
  \task $\res{-12}$
  \task $\res{-99}$
\end{tasks}

\ex{2} \begin{tasks}(4)
  \task $\res{-3}$
  \task $\res{3}$
  \task $\res{-9}$
  \task $\res{4}$ i $\res{-4}$
\end{tasks}

\ex{3} \begin{tasks}(1)
  \task \res{\text{V}}: els negatius estan a l'esquerra del zero.
  \task \res{\text{F}}: $-20$ està més a l'esquerra, $-20<-2$.
  \task \res{\text{V}}: $-2,\ {-1},\ 0,\ 1,\ 2$.
  \task \res{\text{V}}: és el negatiu més pròxim al zero.
  \task \res{\text{V}}: canviar el signe és com mirar la recta en un espill: per exemple, $2<5$ i $-2>-5$.
\end{tasks}


% =====================================================
\seccio{3}{Valor absolut i oposat}

\ex{1} \begin{tasks}(3)
  \task $\lvert -5\rvert =\res{5}$
  \task $\lvert -1\rvert =\res{1}$
  \task $\lvert 0\rvert =\res{0}$
  \task $\lvert 8\rvert =\res{8}$
  \task $\lvert -23\rvert =\res{23}$
  \task $\lvert 7\rvert =\res{7}$
\end{tasks}

\ex{1} \begin{tasks}(3)
  \task $\text{op}(-8)=\res{8}$
  \task $\text{op}(-3)=\res{3}$
  \task $\text{op}(0)=\res{0}$
  \task $\text{op}(4)=\res{-4}$
  \task $\text{op}(-100)=\res{100}$
  \task $\text{op}(15)=\res{-15}$
\end{tasks}

\ex{1} Cadascun està a la meitat de la distància del zero: $36:2=18$. \ Són $\res{18}$ i $\res{-18}$.

\ex{2} $\text{op}(4)=-4$, \ $\lvert -6\rvert =6$, \ $\text{op}(5)=-5$, \ $\lvert 7\rvert =7$, \ $\lvert -2\rvert =2$: \[\res{\text{op}(5)<\text{op}(4)<-2<0<1<\lvert -2\rvert <3<5<\lvert -6\rvert <\lvert 7\rvert }\]

\ex{2} \begin{tasks}(3)
  \task $\lvert -3\rvert \res{=}3$
  \task $\text{op}(-2)\res{>}-1$
  \task $\lvert -2\rvert \res{>}\lvert -1\rvert $
  \task $\text{op}(\text{op}(-1))\res{<}0$
  \task $\lvert -4\rvert \res{>}\text{op}(2)$
  \task $-8\res{=}\text{op}(8)$
\end{tasks}\par (Valors: $3=3$; \ $2>-1$; \ $2>1$; \ $-1<0$; \ $4>-2$; \ $-8=-8$.)

\ex{3} \begin{tasks}(3)
  \task $\res{-7}$
  \task $\res{10}$ o $\res{-10}$
  \task $\res{2}$ o $\res{-2}$
  \task $\res{8}$
  \task $\res{-4}$
  \task Cap: $\text{op}(5)=-5$ i un valor absolut no pot ser negatiu.
\end{tasks}

\ex{3} $\res{-3,\ {-2},\ {-1},\ 0,\ 1,\ 2,\ 3}$. \qquad $\res{-2,\ {-1},\ 0,\ 1,\ 2}$.


% =====================================================
\seccio{4}{Suma i resta d'enters}

\ex{1} \begin{tasks}(5)
  \task $5-7=\res{-2}$
  \task $-8-9=\res{-17}$
  \task $-3+7=\res{4}$
  \task $4-3=\res{1}$
  \task $-5+1=\res{-4}$
  \task $10-14=\res{-4}$
  \task $-7+5=\res{-2}$
  \task $9-13=\res{-4}$
  \task $3+10=\res{13}$
  \task $9-4=\res{5}$
  \task $-5-5=\res{-10}$
  \task $-4+4=\res{0}$
  \task $-3-2=\res{-5}$
  \task $8+12=\res{20}$
  \task $-5+15=\res{10}$
\end{tasks}

\ex{1} \begin{tasks}(3)
  \task $10-3-5=10-8=\res{2}$
  \task $15-9-6=15-15=\res{0}$
  \task $5-8+4=9-8=\res{1}$
  \task $9-3+5=14-3=\res{11}$
  \task $-2+2+7=9-2=\res{7}$
  \task $-10+8+6=14-10=\res{4}$
  \task $-10-3+8=8-13=\res{-5}$
  \task $-4-3-2=\res{-9}$
  \task $-1-5-7=\res{-13}$
  \task $9-2-3=9-5=\res{4}$
  \task $12-4-6=12-10=\res{2}$
  \task $3-7+4=7-7=\res{0}$
  \task $5-9+8=13-9=\res{4}$
  \task $-13+6+4=10-13=\res{-3}$
  \task $-2+10-15=10-17=\res{-7}$
  \task $-11-4+8=8-15=\res{-7}$
  \task $-5-3-4=\res{-12}$
  \task $-8+5+6=11-8=\res{3}$
\end{tasks}

\ex{2} \begin{tasks}(3)
  \task $5+7-2-4=12-6=\res{6}$
  \task $2-6+4-9=6-15=\res{-9}$
  \task $9-6-7+2=11-13=\res{-2}$
  \task $-4-5+3+8=11-9=\res{2}$
  \task $-8+2-7+6=8-15=\res{-7}$
  \task $-1+5+6-7=11-8=\res{3}$
  \task $4-1-2-11=4-14=\res{-10}$
  \task $-2+6-8-1=6-11=\res{-5}$
  \task $-1+2+17-12=19-13=\res{6}$
  \task $-24+15+13-8=28-32=\res{-4}$
  \task $-18-10+12+1=13-28=\res{-15}$
  \task $11-15-13-2=11-30=\res{-19}$
\end{tasks}

\ex{2} \begin{tasks}(1)
  \task $12-4-6+9-15=8-6+9-15=2+9-15=11-15=\res{-4}$; \quad $(12+9)-(4+6+15)=21-25=\res{-4}$
  \task $-7+3-8+10-2+5=-4-8+10-2+5=-12+10-2+5=-2-2+5=-4+5=\res{1}$; \quad $(3+10+5)-(7+8+2)=18-17=\res{1}$
\end{tasks}

\ex{3} \begin{tasks}(3)
  \task $\square=\res{8}$
  \task $\square=\res{5}$
  \task $\square=\res{-10}$
  \task $\square=\res{-9}$
  \task $\square=\res{-5}$
  \task $\square=\res{-3}$
\end{tasks}


% =====================================================
\seccio{5}{Parèntesis i claudàtors}

\ex{1} \begin{tasks}(3)
  \task $(-6)+(-3)=-6-3=\res{-9}$
  \task $(+7)+(-4)=7-4=\res{3}$
  \task $(-8)+(+1)=-8+1=\res{-7}$
  \task $(-3)+(-7)=-3-7=\res{-10}$
  \task $(+12)+(+5)=12+5=\res{17}$
  \task $(+3)+(-5)=3-5=\res{-2}$
  \task $(+3)+(+4)=3+4=\res{7}$
  \task $(+4)+(-6)=4-6=\res{-2}$
  \task $(+7)+(-8)=7-8=\res{-1}$
  \task $(-5)+(+4)=-5+4=\res{-1}$
  \task $(-9)+(+3)=-9+3=\res{-6}$
  \task $(+1)+(-1)=1-1=\res{0}$
  \task $(-7)+(+7)=-7+7=\res{0}$
  \task $(-12)+(-3)=-12-3=\res{-15}$
  \task $(+7)+(+5)=7+5=\res{12}$
  \task $(-4)+(-10)=-4-10=\res{-14}$
  \task $(+6)+(-12)=6-12=\res{-6}$
  \task $(-9)+(+15)=-9+15=\res{6}$
\end{tasks}

\ex{2} \begin{tasks}(1)
  \task $3-(-1)=3+1=\res{4}$
  \task $-8+(-17)=-8-17=\res{-25}$
  \task $-4-(-7)=-4+7=\res{3}$
  \task $-(-5)+3=5+3=\res{8}$
  \task $8+(-14)=8-14=\res{-6}$
  \task $(-2)-(+3)+(-6)=-2-3-6=\res{-11}$
  \task $-(-5)+(-4)-(-10)=5-4+10=\res{11}$
  \task $-15-(-3)+7=-15+3+7=\res{-5}$
  \task $(-8)+(-3)-(-17)+(+1)=-8-3+17+1=\res{7}$
  \task $-2-(-3)+7-8=-2+3+7-8=\res{0}$
  \task $-(-3)+(-1)-(+3)-(-8)=3-1-3+8=\res{7}$
  \task $-3+(-6)-4+1=-3-6-4+1=\res{-12}$
  \task $12-25+3-8-(-6)=12-25+3-8+6=\res{-12}$
  \task $-7+4-(-8)+(-3)=-7+4+8-3=\res{2}$
  \task $-(+4)+(-8)-(-6)+4-1=-4-8+6+4-1=\res{-3}$
  \task $-8+(-3)+4-7-5=-8-3+4-7-5=\res{-19}$
  \task $-(-6)+(-1)-7+(-3)=6-1-7-3=\res{-5}$
  \task $15-(+3)-6+3-12=15-3-6+3-12=\res{-3}$
\end{tasks}

\ex{2} \begin{tasks}(2)
  \task $12+(+3-5)=12+(-2)=\res{10}$
  \task $14+(+12-10)=14+2=\res{16}$
  \task $6-(9-7)=6-2=\res{4}$
  \task $15-(2-9)=15-(-7)=\res{22}$
  \task $11-(-6+3)=11-(-3)=\res{14}$
  \task $10-(-7-5)=10-(-12)=\res{22}$
  \task $13+(-8+2)=13+(-6)=\res{7}$
  \task $17+(-5-9)=17+(-14)=\res{3}$
  \task $1+(-3+8)=1+5=\res{6}$
  \task $9-(2-1)=9-1=\res{8}$
  \task $-6-(2-7)=-6-(-5)=\res{-1}$
  \task $11+(-1-3)=11+(-4)=\res{7}$
  \task $6+(6+3)=6+9=\res{15}$
  \task $-5-(-2-8)=-5-(-10)=\res{5}$
  \task $3+(-8+7)=3+(-1)=\res{2}$
  \task $15+(-6-8)=15+(-14)=\res{1}$
\end{tasks}

\ex{2} \begin{tasks}(2)
  \task $(7-4)+(9-5)=3+(9-5)=3+4=\res{7}$
  \task $(2+6)+(5-8)=8+(5-8)=8+(-3)=\res{5}$
  \task $(5-9)+(2-12)=-4+(2-12)=-4+(-10)=\res{-14}$
  \task $(7+3)-(5+4)=10-(5+4)=10-9=\res{1}$
  \task $(8-12)-(2-5)=-4-(2-5)=-4-(-3)=\res{-1}$
  \task $(10-7)-(-2-6)=3-(-2-6)=3-(-8)=\res{11}$
  \task $-(8+4)+(5-9)=-12+(5-9)=-12+(-4)=\res{-16}$
  \task $-(6-2)-(7-9)=-4-(7-9)=-4-(-2)=\res{-2}$
  \task $(-8-2)+(-1-5)=-10+(-1-5)=-10+(-6)=\res{-16}$
  \task $(7-1)-(-5+11)=6-(-5+11)=6-6=\res{0}$
  \task $(15-8)+(-6-12)=7+(-6-12)=7+(-18)=\res{-11}$
  \task $(-2+3)-(8-4)=1-(8-4)=1-4=\res{-3}$
  \task $(17-12)-(-7-8)=5-(-7-8)=5-(-15)=\res{20}$
  \task $(-10+3)-(2-9)=-7-(2-9)=-7-(-7)=\res{0}$
  \task $-(6-4)+(8+1)=-2+(8+1)=-2+9=\res{7}$
  \task $-(-7-1)-(2-10)=-(-8)-(2-10)=-(-8)-(-8)=\res{16}$
\end{tasks}

\ex{3} \begin{tasks}(1)
  \task $6+[5+(7+2)]=6+[5+9]=6+14=\res{20}$
  \task $8+[4-(3+5)]=8+[4-8]=8+(-4)=\res{4}$
  \task $10-[6+(2+7)]=10-[6+9]=10-15=\res{-5}$
  \task $15-[2-(6-10)]=15-[2-(-4)]=15-6=\res{9}$
  \task $15-[10-(8+4)]=15-[10-12]=15-(-2)=\res{17}$
  \task $12-[7-(2-10)]=12-[7-(-8)]=12-15=\res{-3}$
  \task $(-6)+[5+(2-12)]=(-6)+[5+(-10)]=(-6)+(-5)=\res{-11}$
  \task $(-7)-[3-(4-9)]=(-7)-[3-(-5)]=(-7)-8=\res{-15}$
  \task $(2-10)+[5-(8+2)]=-8+[5-(8+2)]=-8+[5-10]=-8+(-5)=\res{-13}$
  \task $(12-3)-[1-(2-6)]=9-[1-(2-6)]=9-[1-(-4)]=9-5=\res{4}$
  \task $[9-(+5)]+[7+(-10)]=4+[7+(-10)]=4+(-3)=\res{1}$
  \task $[10-(-2)]-[5-(+12)]=12-[5-(+12)]=12-(-7)=\res{19}$
  \task $[8-(6+4)]-(5-7)=[8-10]-(5-7)=-2-(5-7)=-2-(-2)=\res{0}$
  \task $[1+(6-9)]-(8-12)=[1+(-3)]-(8-12)=-2-(8-12)=-2-(-4)=\res{2}$
  \task $2-[6-(-1-21)]=2-[6-(-22)]=2-28=\res{-26}$
  \task $-1+[-5-(2-13)]=-1+[-5-(-11)]=-1+6=\res{5}$
\end{tasks}


% =====================================================
\seccio{6}{Multiplicació i divisió}

\ex{1} \begin{tasks}(3)
  \task $(-4)\cdot (+7)=\res{-28}$
  \task $(+2)\cdot (+6)=\res{12}$
  \task $(-5)\cdot (-7)=\res{35}$
  \task $(+3)\cdot (-8)=\res{-24}$
  \task $(-9)\cdot (-3)=\res{27}$
  \task $(-6)\cdot (+4)=\res{-24}$
  \task $3\cdot (-2)=\res{-6}$
  \task $4\cdot (+5)=\res{20}$
  \task $8\cdot (-6)=\res{-48}$
  \task $-5\cdot (+3)=\res{-15}$
  \task $-2\cdot (-4)=\res{8}$
  \task $-6\cdot (+3)=\res{-18}$
  \task $(+3)\cdot (-5)\cdot (+2)=\res{-30}$
  \task $(-4)\cdot (-1)\cdot (+6)=\res{24}$
  \task $(-2)\cdot (-7)\cdot (-2)=\res{-28}$
  \task $(+5)\cdot (-4)\cdot (-3)=\res{60}$
  \task $(-8)\cdot 4\cdot (-6)=\res{192}$
  \task $5\cdot 3\cdot (-9)\cdot 4=\res{-540}$
\end{tasks}

\ex{1} \begin{tasks}(2)
  \task $(-8):(+2)=\res{-4}$
  \task $(+20):(-10)=\res{-2}$
  \task $(-12):(-4)=\res{3}$
  \task $(-4):(+2)=\res{-2}$
  \task $(+21):(-7)=\res{-3}$
  \task $(-12):(+6)=\res{-2}$
  \task $(-15):(-3)=\res{5}$
  \task $(+32):(+8)=\res{4}$
  \task $(-36):(+9)=\res{-4}$
  \task $(+42):(-7)=\res{-6}$
  \task $(-48):(-8)=\res{6}$
  \task $(+54):(+6)=\res{9}$
  \task $[(+80):(-8)]:(-5)=-10:(-5)=\res{2}$
  \task $[(-70):(-2)]:(-7)=35:(-7)=\res{-5}$
  \task $(+50):[(-30):(+6)]=(+50):(-5)=\res{-10}$
  \task $(-40):[(+24):(+3)]=(-40):8=\res{-5}$
\end{tasks}

\ex{2} \begin{tasks}(6)
  \task $\res{3}$
  \task $\res{8}$
  \task $\res{-7}$
  \task $\res{4}$
  \task $\res{-5}$
  \task $\res{-11}$
\end{tasks}

\ex{2} \begin{tasks}(2)
  \task $[(+6)\cdot (-4)]:(-3)=-24:(-3)=\res{8}$
  \task $[(-15)\cdot (-2)]:(+6)=30:(+6)=\res{5}$
  \task $(-5)\cdot [(+12):(-3)]=(-5)\cdot (-4)=\res{20}$
  \task $[(-5)\cdot (+12)]:(-3)=-60:(-3)=\res{20}$
  \task $24:(-2)\cdot (-3)=(-12)\cdot (-3)=\res{36}$
  \task $(-40)\cdot [(-10):(-5)]=(-40)\cdot 2=\res{-80}$
  \task $(-70)\cdot (-5):(+2)=350:(+2)=\res{175}$
  \task $(+100):(+4)\cdot (-2)=25\cdot (-2)=\res{-50}$
\end{tasks}

\ex{3} \begin{tasks}(3)
  \task $\res{4}$ i $\res{-3}$
  \task $\res{-3}$ i $\res{-6}$
  \task $\res{4}$ i $\res{-5}$
\end{tasks}

\ex{3} \begin{tasks}(1)
  \task Quatre factors negatius (parell): \res{+}. \ $24$
  \task Tres factors negatius (senar): \res{-}. \ $-120$
  \task Set factors negatius (senar): \res{-}. \ $-1$
\end{tasks}


% =====================================================
\seccio{7}{Operacions combinades}

\ex{1} \begin{tasks}(1)
  \task $(-8)\cdot (+2)+(-5)\cdot (-3)=-16+15=\res{-1}$
  \task $(+40):(-8)-(-30):(+6)=-5-(-5)=\res{0}$
  \task $(-2)\cdot (-9)+(-24):(-3)-(-6)\cdot 4=18+8-(-24)=\res{50}$
  \task $(+27):(-3)-(+3)\cdot (-5)-(-6)\cdot (-2)=-9-(-15)-12=\res{-6}$
  \task $(-3)\cdot [(-2)+(-4)]=(-3)\cdot (-6)=\res{18}$
  \task $(+4)\cdot [(-5)+(+2)]=(+4)\cdot (-3)=\res{-12}$
  \task $(+6):[(+5)-(+7)]=(+6):(-2)=\res{-3}$
  \task $(-20):[(-6)-(-2)]=(-20):(-4)=\res{5}$
\end{tasks}

\ex{2} \begin{tasks}(1)
  \task $[(-8)+(+7)]\cdot (-3)=(-1)\cdot (-3)=\res{3}$
  \task $[(-9)+(-3)]:(+6)=(-12):(+6)=\res{-2}$
  \task $19-(-3)\cdot [5-(+8)]=19-(-3)\cdot (-3)=19-9=\res{10}$
  \task $12+(-5)\cdot [8+(-9)]=12+(-5)\cdot (-1)=12+5=\res{17}$
  \task $12-[13-(-7)]:(-5)=12-20:(-5)=12-(-4)=\res{16}$
  \task $10-(+20):[7+(-3)]=10-20:4=10-5=\res{5}$
  \task $(-2)\cdot (5-7)-(-3)\cdot (8-6)=(-2)\cdot (-2)-(-3)\cdot 2=4+6=\res{10}$
  \task $(9-6)\cdot (-2)+(13+3):(-4)=3\cdot (-2)+16:(-4)=-6-4=\res{-10}$
\end{tasks}

\ex{3} \begin{tasks}(1)
  \task $-2\cdot [3-(-4+6)\cdot (-2)]+15:(-5)=-2\cdot [3-2\cdot (-2)]+15:(-5)=-2\cdot 7+(-3)=-14-3=\res{-17}$
  \task $[(-3)\cdot 4-(-2)\cdot 5]:(-2)=[-12+10]:(-2)=(-2):(-2)=\res{1}$
  \task $18:[-7-(-4)]-(-2)\cdot [5-3\cdot (-1)]=18:(-3)-(-2)\cdot 8=-6+16=\res{10}$
  \task $(-4)\cdot [6-2\cdot (3-7)]-[-10:(-5)]=(-4)\cdot [6+8]-2=-56-2=\res{-58}$
\end{tasks}

\ex{3} \begin{tasks}(3)
  \task $\res{(6-4)\cdot(2-5)}=-6$
  \task $\res{(-3+5)}\cdot2-4=0$
  \task $10-2\cdot\res{(3+1)}=2$
\end{tasks}


% =====================================================
\seccio{8}{Problemes amb enters}

\ex{1} $-18-4=\res{-22}$\,°C (baixa 22 graus). \ $4-(-18)=4+18=\res{+22}$\,°C (puja 22 graus).

\ex{1} $0+5-3+5-8+10+5-6=25-17=\res{8}$: \ està en la planta 8.

\ex{1} $-14+9=-5$: \ encara li deu $\res{5\text{ €}}$.

\ex{1} $-14+10-5=\res{-9}$\,°C.

\ex{1} $8-(-2)=8+2=\res{10\text{ m}}$.

\ex{1} $-5+3+8-2=11-7=\res{+4}$: \ ha pujat 4 cèntims.

\ex{2} Al migdia: $-3+9=6$\,°C. \ Diferència: $-1-6=\res{-7}$\,°C: \ ha baixat 7 graus.

\ex{2} $-265-60=-325$: \ va nàixer l'any $\res{325\text{ a.\,C.}}$

\ex{2} $3478-(-1441)=3478+1441=\res{4919\text{ m}}$.

\ex{2} $850+5\cdot250-2300=850+1250-2300=\res{-200\text{ €}}$: \ el compte està en descobert (deu 200\,€).

\ex{3} $-212-(-287)=-212+287=\res{75\text{ anys}}$.

\ex{3} $-120+45-3\cdot45=-120+45-135=\res{-210}$: \ a 210\,m de profunditat.

\ex{3} $11\cdot3+6\cdot(-2)+3\cdot0=33-12=\res{21\text{ punts}}$. \ Fallar-ne 9 resta $18$: cal sumar-ne almenys 18, és a dir, $18:3=\res{6}$ encerts (en queden $20-9-3=8$, així que és possible).


\end{document}
